% Part: first-order-logic
% Chapter: completeness
% Section: identity

\documentclass[../../../include/open-logic-section]{subfiles}

\begin{document}

\olfileid{fol}{com}{ide}
\olsection{الهوية}

\begin{explain}
أثبتنا بالبناء السابق ما يلزم لاكتمال منطق الرتبة الأولى حين تخلو
المجموعات~${\OLId{greek-variable}{Gamma}}$ من~$\eq$: فإن نموذج الحدود يرضي كل
$!A \in {\OLId{greek-variable}{Gamma}}^*$ خالية من~$\eq$، فيرضي كل $!A \in {\OLId{greek-variable}{Gamma}}$.
وأما مع وجود $\eq$ فلا يفي ذلك البناء بالغرض. فقد تضم ${\OLId{greek-variable}{Gamma}}^*$
!!a{sentence}~$\eq[{\OLId{latin-ordinary}{t}}][{\OLId{latin-ordinary}{t}}']$، وقيمة الحد في نموذج الحدود هي نفسه.
فإذا اختلف الحدان ${\OLId{latin-ordinary}{t}}$ و${\OLId{latin-ordinary}{t}}'$، اختلفت قيمتاهما، وهما ${\OLId{latin-ordinary}{t}}$ و${\OLId{latin-ordinary}{t}}'$
على الترتيب، فتكذب $\eq[{\OLId{latin-ordinary}{t}}][{\OLId{latin-ordinary}{t}}']$. وعلاج هذا الخلل ببناء يسمى «أخذ
خارج القسمة».
\end{explain}

\begin{defn}
  إذا كانت ${\OLId{greek-variable}{Gamma}}^*$ مجموعة جمل متسقة و!!{complete} في~$\Lang L$،
  وضعنا العلاقة $\approx$ بين الحدود المغلقة في~$\Lang L$ على هذا الوجه:
  \[
  {\OLId{latin-ordinary}{t}} \approx {\OLId{latin-ordinary}{t}}' \text{\quad إذا وفقط إذا \quad} \eq[{\OLId{latin-ordinary}{t}}][{\OLId{latin-ordinary}{t}}'] \in {\OLId{greek-variable}{Gamma}}^*
  \]
\end{defn}

\begin{prop}
\ollabel{prop:approx-equiv}
للعلاقة $\approx$ الخصائص الآتية:
\begin{enumerate}
\item $\approx$ انعكاسية.
\item $\approx$ تماثلية.
\item  $\approx$ تعدية.
\item إذا كان ${\OLId{latin-ordinary}{t}} \approx {\OLId{latin-ordinary}{t}}'$، وكانت ${\OLId{latin-ordinary}{f}}$ !!a{function}، وكانت ${\OLId{latin-ordinary}{t}}_{\OLMathNumeral{1}}$،
  \dots، ${\OLId{latin-ordinary}{t}}_{{\OLId{latin-ordinary}{i}}-{\OLMathNumeral{1}}}$، ${\OLId{latin-ordinary}{t}}_{{\OLId{latin-ordinary}{i}}+{\OLMathNumeral{1}}}$، \dots، ${\OLId{latin-ordinary}{t}}_{\OLId{latin-ordinary}{n}}$ حدودًا مغلقة، فإن
\[
\Atom{f}{t_1,\dots, t_{i-1}, t, t_{i+1}, \dots, t_n} \approx
\Atom{f}{t_1,\dots, t_{i-1}, t', t_{i+1}, \dots, t_n}.
\]
\item إذا كان ${\OLId{latin-ordinary}{t}} \approx {\OLId{latin-ordinary}{t}}'$، وكان ${\OLId{latin-looped}{R}}$ !!a{predicate}، وكانت ${\OLId{latin-ordinary}{t}}_{\OLMathNumeral{1}}$،
  \dots، ${\OLId{latin-ordinary}{t}}_{{\OLId{latin-ordinary}{i}}-{\OLMathNumeral{1}}}$، ${\OLId{latin-ordinary}{t}}_{{\OLId{latin-ordinary}{i}}+{\OLMathNumeral{1}}}$، \dots، ${\OLId{latin-ordinary}{t}}_{\OLId{latin-ordinary}{n}}$ حدودًا مغلقة، فإن
\begin{multline*}
\Atom{R}{t_1,\dots, t_{i-1}, t, t_{i+1}, \dots, t_n} \in {\OLId{greek-variable}{Gamma}}^* \text{ إذا وفقط إذا } \\
\Atom{R}{t_1,\dots, t_{i-1}, t', t_{i+1}, \dots, t_n} \in {\OLId{greek-variable}{Gamma}}^*.
\end{multline*}
\end{enumerate}
\end{prop}

\begin{proof}
الاتساق وكون ${\OLId{greek-variable}{Gamma}}^*$ !!{complete} يجعلان $\eq[{\OLId{latin-ordinary}{t}}][{\OLId{latin-ordinary}{t}}'] \in
{\OLId{greek-variable}{Gamma}}^*$ مكافئًا لـ${\OLId{greek-variable}{Gamma}}^* \Proves \eq[{\OLId{latin-ordinary}{t}}][{\OLId{latin-ordinary}{t}}']$؛ فتنحصر مطالب
البرهان في الأحكام الآتية:
\begin{enumerate}
\item ${\OLId{greek-variable}{Gamma}}^* \Proves \eq[{\OLId{latin-ordinary}{t}}][{\OLId{latin-ordinary}{t}}]$ لكل الحدود المغلقة~${\OLId{latin-ordinary}{t}}$.
\item إذا كان ${\OLId{greek-variable}{Gamma}}^* \Proves \eq[{\OLId{latin-ordinary}{t}}][{\OLId{latin-ordinary}{t}}']$، فإن
  ${\OLId{greek-variable}{Gamma}}^* \Proves \eq[{\OLId{latin-ordinary}{t}}'][{\OLId{latin-ordinary}{t}}]$.
\item إذا كان ${\OLId{greek-variable}{Gamma}}^* \Proves \eq[{\OLId{latin-ordinary}{t}}][{\OLId{latin-ordinary}{t}}']$ و${\OLId{greek-variable}{Gamma}}^* \Proves
  \eq[{\OLId{latin-ordinary}{t}}'][{\OLId{latin-ordinary}{t}}'']$، فإن ${\OLId{greek-variable}{Gamma}}^* \Proves \eq[{\OLId{latin-ordinary}{t}}][{\OLId{latin-ordinary}{t}}'']$.
\item إذا كان ${\OLId{greek-variable}{Gamma}}^* \Proves \eq[{\OLId{latin-ordinary}{t}}][{\OLId{latin-ordinary}{t}}']$، فإن
\[
{\OLId{greek-variable}{Gamma}}^* \Proves
\eq[\Atom{f}{t_1,\dots,t_{i-1},t,t_{i+1},\dots,t_n}][\Atom{f}{t_1,\dots,t_{i-1},t',t_{i+1},\dots,t_n}]
\]
لكل !!{function}~${\OLId{latin-ordinary}{f}}$ ذات ${\OLId{latin-ordinary}{n}}$ مواضع ولكل حدود مغلقة ${\OLId{latin-ordinary}{t}}_{\OLMathNumeral{1}}$، \dots،
${\OLId{latin-ordinary}{t}}_{{\OLId{latin-ordinary}{i}}-{\OLMathNumeral{1}}}$، ${\OLId{latin-ordinary}{t}}_{{\OLId{latin-ordinary}{i}}+{\OLMathNumeral{1}}}$، \dots،~${\OLId{latin-ordinary}{t}}_{\OLId{latin-ordinary}{n}}$.
\item إذا كان ${\OLId{greek-variable}{Gamma}}^* \Proves \eq[{\OLId{latin-ordinary}{t}}][{\OLId{latin-ordinary}{t}}']$ و
${\OLId{greek-variable}{Gamma}}^* \Proves
\Atom{R}{t_1,\dots,t_{i-1},t,t_{i+1},\dots,t_n}$، فإن
${\OLId{greek-variable}{Gamma}}^* \Proves \Atom{R}{t_1,\dots,t_{i-1},t',t_{i+1},\dots,t_n}$
لكل !!{predicate}~${\OLId{latin-looped}{R}}$ ذي ${\OLId{latin-ordinary}{n}}$ مواضع ولكل حدود مغلقة ${\OLId{latin-ordinary}{t}}_{\OLMathNumeral{1}}$، \dots،
${\OLId{latin-ordinary}{t}}_{{\OLId{latin-ordinary}{i}}-{\OLMathNumeral{1}}}$، ${\OLId{latin-ordinary}{t}}_{{\OLId{latin-ordinary}{i}}+{\OLMathNumeral{1}}}$، \dots،~${\OLId{latin-ordinary}{t}}_{\OLId{latin-ordinary}{n}}$.
\end{enumerate}
\end{proof}

\begin{prob}
أكمل برهان~\olref[fol][com][ide]{prop:approx-equiv}.
\end{prob}

\begin{defn}
لتكن ${\OLId{greek-variable}{Gamma}}^*$ مجموعة متسقة و!!{complete} في لغة~$\Lang L$، وليكن
${\OLId{latin-ordinary}{t}}$ حدًا مغلقًا، والعلاقة $\approx$ هي المتقدمة. نرمز بـ$\Trm[{\OLId{latin-looped}{L}}]_{\OLMathNumeral{0}}$
إلى مجموعة الحدود المغلقة في~$\Lang L$، ونضع:
\[
\equivrep{{\OLId{latin-ordinary}{t}}}{\approx} = \Setabs{{\OLId{latin-ordinary}{t}}'}{{\OLId{latin-ordinary}{t}}'\in \Trm[{\OLId{latin-looped}{L}}]_{\OLMathNumeral{0}}, {\OLId{latin-ordinary}{t}} \approx {\OLId{latin-ordinary}{t}}'}
\]
و$\equivclass{\Trm[{\OLId{latin-looped}{L}}]_{\OLMathNumeral{0}}}{\approx} = \Setabs{\equivrep{{\OLId{latin-ordinary}{t}}}{\approx}}{{\OLId{latin-ordinary}{t}} \in \Trm[{\OLId{latin-looped}{L}}]_{\OLMathNumeral{0}}}$.
\end{defn}

\begin{defn}
\ollabel{defn:term-model-factor}
خذ $\Struct M = \Struct M({\OLId{greek-variable}{Gamma}}^*)$، نموذج الحدود الخاص بـ~${\OLId{greek-variable}{Gamma}}^*$
في \olref[mod]{defn:termmodel}. ونعرّف
$\Struct{\equivclass{M}{\approx}}$ بأنها ال!!{structure} المحددة بما يأتي:
\begin{enumerate}
\item $\Domain{\equivclass{{\OLId{latin-looped}{M}}}{\approx}} = \equivclass{\Trm[{\OLId{latin-looped}{L}}]_{\OLMathNumeral{0}}}{\approx}$.
\item $\Assign{{\OLId{latin-ordinary}{c}}}{\equivclass{{\OLId{latin-looped}{M}}}{\approx}} = \equivrep{{\OLId{latin-ordinary}{c}}}{\approx}$
\item $\Assign{{\OLId{latin-ordinary}{f}}}{\equivclass{{\OLId{latin-looped}{M}}}{\approx}}(\equivrep{{\OLId{latin-ordinary}{t}}_{\OLMathNumeral{1}}}{\approx}, \dots,
  \equivrep{{\OLId{latin-ordinary}{t}}_{\OLId{latin-ordinary}{n}}}{\approx}) = \equivrep{\Atom{f}{t_1,\dots, t_n}}{\approx}$
\item $\tuple{\equivrep{{\OLId{latin-ordinary}{t}}_{\OLMathNumeral{1}}}{\approx}, \dots, \equivrep{{\OLId{latin-ordinary}{t}}_{\OLId{latin-ordinary}{n}}}{\approx}} \in
  \Assign{{\OLId{latin-looped}{R}}}{\equivclass{{\OLId{latin-looped}{M}}}{\approx}}$ إذا وفقط إذا كان
  $\Sat{{\OLId{latin-looped}{M}}}{\Atom{R}{t_1,\dots, t_n}}$؛ أي إذا وفقط إذا كان
  $\Atom{R}{t_1,\dots, t_n} \in {\OLId{greek-variable}{Gamma}}^*$.
\end{enumerate}
\end{defn}

\begin{explain}
في تعريف كل من $\Assign{{\OLId{latin-ordinary}{f}}}{\equivclass{{\OLId{latin-looped}{M}}}{\approx}}$ و
$\Assign{{\OLId{latin-looped}{R}}}{\equivclass{{\OLId{latin-looped}{M}}}{\approx}}$ سمّينا عناصر
$\equivclass{\Trm[{\OLId{latin-looped}{L}}]_{\OLMathNumeral{0}}}{\approx}$ بأسماء من الصورة
$\equivrep{{\OLId{latin-ordinary}{t}}}{\approx}$؛ فاعتمدنا في وصفها على \emph{ممثلين}~${\OLId{latin-ordinary}{t}} \in
\equivrep{{\OLId{latin-ordinary}{t}}}{\approx}$. ولا يصح التعريف إلا إذا استغنى ناتجه عن اختيار
الممثل: فلو اخترنا~${\OLId{latin-ordinary}{t}}'$ مع بقاء الصنف نفسه
($\equivrep{{\OLId{latin-ordinary}{t}}}{\approx} = \equivrep{{\OLId{latin-ordinary}{t}}'}{\approx}$)، وجب ألا تتغير
النتيجة. انظر مثلًا إلى ${\OLId{latin-looped}{R}}$ وهو !!{predicate} أحادي الموضع.
يقضي البند الأخير بأن $\equivrep{{\OLId{latin-ordinary}{t}}}{\approx} \in \Assign{{\OLId{latin-looped}{R}}}{\equivclass{{\OLId{latin-looped}{M}}}{\approx}}$
إذا وفقط إذا كان $\Sat{{\OLId{latin-looped}{M}}}{\Atom{R}{t}}$. فلو أخذنا حدًا آخر~${\OLId{latin-ordinary}{t}}'$
% تصحيح مصدر (0133-I1): صحّحنا حدَّ الفرض من t إلى الممثل البديل t'، وهو خطأ صوري موروث من الأصل وخط الأساس.
حيث ${\OLId{latin-ordinary}{t}} \approx {\OLId{latin-ordinary}{t}}'$، وافترضنا $\Sat/{{\OLId{latin-looped}{M}}}{\Atom{R}{t'}}$، لاقتضى
التعريف $\equivrep{{\OLId{latin-ordinary}{t}}'}{\approx} \notin \Assign{{\OLId{latin-looped}{R}}}{\equivclass{{\OLId{latin-looped}{M}}}{\approx}}$.
لكن ${\OLId{latin-ordinary}{t}} \approx {\OLId{latin-ordinary}{t}}'$ يوجب $\equivrep{{\OLId{latin-ordinary}{t}}}{\approx} =
\equivrep{{\OLId{latin-ordinary}{t}}'}{\approx}$، ويستحيل الجمع بين
$\equivrep{{\OLId{latin-ordinary}{t}}}{\approx} \in \Assign{{\OLId{latin-looped}{R}}}{\equivclass{{\OLId{latin-looped}{M}}}{\approx}}$ و
$\equivrep{{\OLId{latin-ordinary}{t}}}{\approx} \notin \Assign{{\OLId{latin-looped}{R}}}{\equivclass{{\OLId{latin-looped}{M}}}{\approx}}$.
وتمنع \olref{prop:approx-equiv} وقوع هذا الاختلاف.
\end{explain}

\begin{prop}
$\Struct{\equivclass{M}{\approx}}$ معرّفة تعريفًا سليمًا؛ أي إذا كانت
${\OLId{latin-ordinary}{t}}_{\OLMathNumeral{1}}$، \dots، ${\OLId{latin-ordinary}{t}}_{\OLId{latin-ordinary}{n}}$، ${\OLId{latin-ordinary}{t}}_{\OLMathNumeral{1}}'$، \dots، ${\OLId{latin-ordinary}{t}}_{\OLId{latin-ordinary}{n}}'$ حدودًا مغلقة، وكان
${\OLId{latin-ordinary}{t}}_{\OLId{latin-ordinary}{i}} \approx {\OLId{latin-ordinary}{t}}_{\OLId{latin-ordinary}{i}}'$ لكل ${\OLId{latin-ordinary}{i}}$ يحقق ${\OLMathNumeral{1}} \le {\OLId{latin-ordinary}{i}} \le {\OLId{latin-ordinary}{n}}$، فإن
\begin{enumerate}
\item $\equivrep{\Atom{f}{t_1,\dots, t_n}}{\approx} =
    \equivrep{\Atom{f}{t_1',\dots, t_n'}}{\approx}$؛ أي إن
  \[
  \Atom{f}{t_1,\dots, t_n} \approx \Atom{f}{t_1',\dots, t_n'}
  \]
  و
\item $\Sat{{\OLId{latin-looped}{M}}}{\Atom{R}{t_1,\dots, t_n}}$ إذا وفقط إذا كان
  $\Sat{{\OLId{latin-looped}{M}}}{\Atom{R}{t_1',\dots, t_n'}}$؛ أي إن
  \[
    \Atom{R}{t_1,\dots, t_n} \in {\OLId{greek-variable}{Gamma}}^* \text{ إذا وفقط إذا }
    \Atom{R}{t_1',\dots, t_n'} \in {\OLId{greek-variable}{Gamma}}^*.
  \]
\end{enumerate}
\end{prop}

\begin{proof}
نطبق \olref{prop:approx-equiv} استقراءً على~${\OLId{latin-ordinary}{n}}$، فيثبت المطلوب.
\end{proof}

وقبل إثبات لمّة الصدق نحتاج إلى لمّة القيم، كما احتجنا إليها في نموذج الحدود.

\begin{lem}
\ollabel{lem:val-in-termmodel-factored} ليكن $\Struct M = \Struct M
 ({\OLId{greek-variable}{Gamma}}^*)$؛ عندئذ $\Value{{\OLId{latin-ordinary}{t}}}{\equivclass{{\OLId{latin-looped}{M}}}{\approx}} = \equivrep{{\OLId{latin-ordinary}{t}}}
 {\approx}$ لكل حد مغلق~${\OLId{latin-ordinary}{t}}$.
\end{lem}
\begin{proof}
يسلك البرهان مسلك \olref[mod]{lem:val-in-termmodel} نفسه.
\end{proof}

\begin{prob}
أكمل برهان~\olref[fol][com][ide]{lem:val-in-termmodel-factored}.
\end{prob}

\begin{lem}
\ollabel{lem:truth}
$\Sat{\equivclass{{\OLId{latin-looped}{M}}}{\approx}}{!A}$ إذا وفقط إذا كان
$!A \in {\OLId{greek-variable}{Gamma}}^*$ لكل الجمل~$!A$.
\end{lem}

\begin{proof}
نستقرئ~$!A$ على نحو برهان \olref[mod]{lem:truth}. ولا يبقى من
الحالات ما يستدعي زيادة بيان إلا $!A \ident \eq[{\OLId{latin-ordinary}{t}}][{\OLId{latin-ordinary}{t}}']$، وبيانها:
\begin{align*}
\Sat{\equivclass{{\OLId{latin-looped}{M}}}{\approx}}{\eq[{\OLId{latin-ordinary}{t}}][{\OLId{latin-ordinary}{t}}']} & \text{ إذا وفقط إذا } \equivrep{{\OLId{latin-ordinary}{t}}}{\approx} = \equivrep{{\OLId{latin-ordinary}{t}}'}{\approx}
\text{ (بحسب تعريف $\Struct{\equivclass{M}{\approx}}$)}\\
& \text{ إذا وفقط إذا } {\OLId{latin-ordinary}{t}} \approx {\OLId{latin-ordinary}{t}}' \text{ (بحسب تعريف $\equivrep{{\OLId{latin-ordinary}{t}}}{\approx}$)}\\
& \text{ إذا وفقط إذا } \eq[{\OLId{latin-ordinary}{t}}][{\OLId{latin-ordinary}{t}}'] \in {\OLId{greek-variable}{Gamma}}^* \text{ (بحسب تعريف $\approx$).}
\end{align*}
\end{proof}

\begin{digress}
إن $\Struct{M(\Gamma^*)}$ !!{enumerable} ولانهائي دائمًا، لكن
$\Struct{\equivclass{M}{\approx}}$ قد يكون منتهيًا؛ إذ ربما لم يبق
إلا عدد منتهٍ من الأصناف $\equivrep{{\OLId{latin-ordinary}{t}}}{\approx}$. ولا غرابة في
ذلك، فقد تتضمن ${\OLId{greek-variable}{Gamma}}$ !!{sentence}s لا تصدق إلا في
!!{structure} منتهية. ومثال ذلك $\lforall[{\OLId{latin-ordinary}{x}}][\lforall[{\OLId{latin-ordinary}{y}}][{\OLId{latin-ordinary}{x}} =
{\OLId{latin-ordinary}{y}}]]$؛ فهي !!{sentence} متسقة، ولا ترضيها إلا !!{structure}s
التي لا يوجد في !!a{domain} لها إلا !!{element} واحد.
\end{digress}

\end{document}
